\begin{table}[!tbp]\centering
\footnotesize
\caption{Fiscal magnitudes in 2023, million lei: micro-enterprise tax paid against 16 percent of reported pre-tax profit, by revenue band}\label{tab:fiscal}
\resizebox{\linewidth}{!}{%
\begin{tabular}{lrrrrrrrr}\toprule
Revenue band (\euro{}, prev.\ rate) & Firms & Micro firms & Micro tax paid & 16\% of their profit & Gap & Gap share (\%) & Profit-tax firms & Profit tax paid \\ \midrule
< 60k & 420,646 & 210,267 & 306 & 1,187 & 880 & 12 & 196,922 & 529 \\
60-250k & 164,407 & 143,787 & 822 & 3,486 & 2,664 & 36 & 19,880 & 335 \\
250-500k & 55,483 & 48,765 & 785 & 3,217 & 2,432 & 33 & 6,527 & 176 \\
500k-1M & 31,385 & 21,463 & 927 & 1,819 & 891 & 12 & 9,858 & 471 \\
> 1M & 44,797 & 3,762 & 502 & 982 & 480 & 6 & 40,729 & 20,104 \\
\bottomrule\end{tabular}}
\begin{minipage}{0.95\linewidth}\vspace{2pt}\footnotesize Micro firms are those registered for the micro-enterprise tax in the ANAF registry of 31 August 2023, profit-tax firms those registered for profit tax; firms above \euro{}500,000 registered as micro-enterprises are in-year crossers, taxed on revenue until the quarter of crossing and on profit after it, and their tax paid includes both, as are hospitality firms, which faced no ceiling in 2023. Profit at 16\% uses reported pre-tax profit floored at zero and is an upper bound, since Section~\ref{sec:base} shows reported margins fall by 8 to 17 percent of their level when firms move to profit tax. The gap is the static revenue difference; the share column is each band's share of the total gap.\end{minipage}
\end{table}